% TYK §2.1.9 + Ek-8: alfabetik sıra; yabancı dildeki kısaltmaların Türkçe
% karşılığı **parantez içinde** verilir.
\begin{simge}
	\item [AI]     Artificial Intelligence (Yapay Zekâ)
	\item [AIOps]  Artificial Intelligence for IT Operations (BT Operasyonları için Yapay Zekâ)
	\item [API]    Application Programming Interface (Uygulama Programlama Arayüzü)
	\item [AUC]    Area Under the Curve (Eğri Altındaki Alan)
	\item [FN]     False Negative (Yanlış Negatif)
	\item [FP]     False Positive (Yanlış Pozitif)
	\item [GBT]    Gradient Boosted Trees (Gradyan Artırmalı Ağaçlar)
	\item [IaaS]   Infrastructure as a Service (Hizmet Olarak Altyapı)
	\item [IQR]    Interquartile Range (Çeyrekler Arası Aralık)
	\item [JSON]   JavaScript Object Notation
	\item [LGBM]   Light Gradient Boosting Machine (Hafif Gradyan Artırma Makinesi)
	\item [LLM]    Large Language Model (Büyük Dil Modeli)
	\item [LR]     Linear Regression (Doğrusal Regresyon)
	\item [MAE]    Mean Absolute Error (Ortalama Mutlak Hata)
	\item [MAPE]   Mean Absolute Percentage Error (Ortalama Mutlak Yüzdelik Hata)
	\item [ML]     Machine Learning (Makine Öğrenmesi)
	\item [MLP]    Multilayer Perceptron (Çok Katmanlı Algılayıcı)
	\item [RMSE]   Root Mean Squared Error (Kök Ortalama Kare Hata)
	\item [ROC]    Receiver Operating Characteristic (Alıcı İşletim Karakteristiği)
	\item [RF]     Random Forest (Rastgele Orman)
	\item [SHAP]   SHapley Additive exPlanations (Shapley Toplamsal Açıklamalar)
	\item [TSDB]   Time Series Database (Zaman Serisi Veritabanı)
	\item [TYK]    Tez Yazım Kılavuzu
	\item [VM]     Virtual Machine (Sanal Makine)
	\item [XAI]    Explainable Artificial Intelligence (Açıklanabilir Yapay Zekâ)
	\item [XGB]    eXtreme Gradient Boosting (Aşırı Gradyan Artırma)
\end{simge}
